\documentclass[11pt]{article}
\usepackage[margin=1in]{geometry}
\usepackage{amsmath,amssymb,amsthm,booktabs,microtype,hyperref}
\hypersetup{hidelinks,pdfauthor={T},pdftitle={Recovery Certificates Under Bounded Provenance Omissions}}
\newtheorem{theorem}{Theorem}
\title{Recovery Certificates Under Bounded Provenance Omissions}
\author{T}
\date{}
\begin{document}
\maketitle
\begin{abstract}
An operator restoring an agent needs to know whether an incomplete
dependency log could hide a path from compromised state to the
restored service. We measure the smallest omission that invalidates
a restoration certificate. Counting omitted edges is exact for an
edge budget and can be arbitrarily optimistic when an entire log
source can disappear. We give an exact source-based algorithm and
an additive certificate that can sometimes settle the recovery decision
without enumerating source combinations. On acyclic graphs, an exact
algorithm stores only labels that can recur after the current vertex.
The maximum number of such labels controls its state space. Separate
source-subset enumeration checks both exact algorithms on synthetic graphs.
An intervention can certify a logging source or rebuild an artifact.
The two choices share a path-hitting condition, and optimizing their
cost can be hard even when detecting a suspect path is simple.
\end{abstract}

\section*{Summary and scope}
\paragraph{Question.} How many missing logging sources could conceal a
path from compromised state to an artifact being restored?
\paragraph{Finding.} The source-omission radius is the smallest number of
distinct source labels along a suspect path. The exact search can
forget labels that cannot recur, and an additive certificate can rule
out a path under the declared source budget.
\paragraph{Evidence and scope.} The result is proved and tested
against source-subset enumeration. An incident team can use the
certificate to decide whether it needs more evidence before restoring an
artifact. The guarantee requires a complete candidate dependency graph.

\section{The log that went quiet}
Suppose an incident team rebuilds an agent's tool adapter and wants
to restore a previously saved plan. The plan may still depend on
compromised state through cached outputs. A missing log is especially
important when it records several steps of one dependency chain.
Counting those steps as several unrelated omissions gives a
misleading picture of how much evidence must fail.

Let $G=(V,E_0\cup E_*)$ be a finite directed candidate dependency
graph after removing independently rebuilt artifacts from possible
taint paths. The source $s$ represents suspect state and the target
$t$ is an artifact whose recovery is proposed. Edges in $E_0$ are
included in all modeled histories. Edges in $E_*$ are possible
dependencies absent from the observed log. Each such edge has one
source label $\ell(e)$. Activating a source permits all candidate
edges with its label. We include all subsets of those permitted
edges as possibilities. This models omission of an entire logging
channel rather than loss of individual records.

The graph and its labels require validation outside this analysis.
If a feasible dependency is absent even from $E_*$, or if a
rebuild does not exclude suspect inputs, the model can miss a
taint path. A certificate records the candidate graph, the
assumed maximum number of affected sources, and the conditions
that would invalidate either. Logs can justify an upper omission
budget only through evidence external to the missing entries.

\section{The minimum omission that opens a path}
For an $s$-to-$t$ path $P$ write
\[
m(P)=|P\cap E_*|,\qquad
S(P)=\{\ell(e):e\in P\cap E_*\}.
\]
Take a minimum over the empty path family to be infinity.
\begin{theorem}
If any set of at most $k$ optional edges can be absent from the log,
a hidden taint path is possible exactly when
$\min_P m(P)\le k$. If at most $k$ logging sources can be absent,
the corresponding condition is $\min_P |S(P)|\le k$.
\end{theorem}
\begin{proof}
A feasible hidden path uses only optional edges or source labels
admitted by the omission set. Its cost is therefore at most the
budget. Conversely, activate precisely the optional edges or
labels of a path within the budget. The path is now present.
\end{proof}

We call the second minimum the source radius $r$. A known path
gives $r=0$. A proposed restoration is safe against the stipulated
omission budget only if $r>k$. A chain of $L$ optional edges all
recorded by one source has edge radius $L$ and source radius one.
No constant factor can turn the edge radius into a lower
certificate for the source radius.

The edge version is a shortest path with zero or unit costs.
The source version is a minimum-color path with free observed
edges. The difficulty of using few distinct colors is
established in the graph literature \cite{broersma}. In the
directed setting, a direct reduction from set cover puts one
layer boundary per universe element. Each arc crossing a
boundary has a color corresponding to a set that contains that
element. A path selects one such arc per element, and its
distinct colors form a cover. Parallel arcs can be removed by
inserting a private intermediate vertex and a free edge.
Thus a polynomial exact algorithm is not expected in general.

\section{A fast certificate and an exact fallback}
Assume the candidate graph is acyclic, as in a dependency log
ordered by creation. For each optional source $g$ let $M_g$
be the maximum number of its edges on any $s$-to-$t$ path.
Compute $M_g$ by a longest-path pass with weight one on its
edges and zero elsewhere. These passes use $O(g(|V|+|E|))$
operations for $g$ sources, apart from rational arithmetic.
Ignore groups with $M_g=0$.
Give an optional edge labeled $g$ weight $1/M_g$ and every
observed edge weight zero. On every path, each group contributes
at most one, so the path's additive cost is at most $|S(P)|$.
If $L$ is the minimum additive cost and $P_L$ attains it, then
\[
\lceil L\rceil \le r \le |S(P_L)|.
\]
The left inequality follows by minimizing a pointwise bound and
rounding because $r$ is an integer. The right inequality comes
from a concrete path. A candidate graph with no possible path
can be recognized first by reachability.

If $\lceil L\rceil>k$, the bound certifies the restoration in
the stated model. If $|S(P_L)|\le k$, that witness disproves
the guarantee. Otherwise use an exact algorithm. Keep, at
each vertex, inclusion-minimal sets of source labels that reach
it from $s$. Extend a set by taking its union with the label
of the outgoing edge, then discard supersets. Removing a
superset is safe because appending the same future labels
preserves inclusion. Processing vertices in topological order
yields the minimal sets at $t$, and their smallest cardinality
is $r$. At most $2^g$ masks can occur when there are $g$ possible
sources. For a decision at a fixed budget, sets exceeding that
budget can be discarded early.

The certificate makes it possible to separate proof obligations.
The cheap lower bound rules out all admitted hidden paths. The
upper bound supplies a path and the exact sources whose absence
would enable it. A recovery team can then look for independent
evidence about those sources or rebuild an internal artifact on
the witness path, and repeat the check on the amended graph.
The resulting sequence is an auditable search for a sufficient
repair. It does not by itself minimize the total repair cost.

\section{Remembering only sources that can recur}
The full antichain can contain labels that will never influence
a later choice. At a vertex $v$ on some suspect path, let $B(v)$ be
the union of optional labels on all $s$-to-$v$ paths and let $A(v)$
be the union on all $v$-to-$t$ paths. Call
$L(v)=B(v)\cap A(v)$ the live labels. Two passes over a topological
ordering compute these unions, ignoring vertices that cannot lie
on an $s$-to-$t$ path.

\begin{theorem}
For a prefix $P$ ending at $v$, store its full cost $|S(P)|$ and
retain only the mask $S(P)\cap L(v)$. At every vertex, keeping the
lowest cost for each mask preserves the exact source radius.
If $w=\max_v|L(v)|$, the algorithm processes at most
$O(2^w(|V|+|E|))$ states and transitions, with one mask operation
and a table update per transition up to a constant factor.
\end{theorem}
\begin{proof}
Suppose $P$ ends at $u$ with mask $M$. Extend it along $u\to v$.
Let $Q$ be the singleton optional label of that edge, or the empty
set for an observed edge. For any edge on an $s$-to-$t$ path,
$Q\subseteq A(u)$. Consequently $Q\cap S(P)=Q\cap M$, and
the cost increments by $|Q\setminus M|$. Also
$A(v)\subseteq A(u)$, since a suffix from $v$ can follow the edge
from $u$. Any previously used label that belongs to $L(v)$ was
therefore already in $M$. The next mask is
$(M\cup Q)\cap L(v)$. The increment and next mask depend only
on the current mask, so a larger cost with that mask can be
discarded. Start at $s$ with the empty mask and cost zero.
Induction over vertices gives the minimum prefix cost for each
reachable mask. At $t$, the suffix has no edges and the mask
is empty, so its stored cost is the radius. Each table has at
most $2^w$ entries. Iterating over its outgoing edges yields
the stated bound, excluding bit-set and table-operation costs.
\end{proof}

This bound is useful when sources are numerous but each records
a localized portion of a dependency history. Consider a chain of
$k$ diamonds. Within each diamond, its two branches are each
two optional edges of one label. Each diamond uses two labels
that occur nowhere else. At a branch vertex the live width is
one, and at a join it is zero. The full antichain at the target
contains all $2^k$ choices, each with one label from every
diamond. The live-label algorithm stores one cost per vertex,
for $3k+1$ states in total, and returns radius $k$.

\section{Independent checks}
An antichain solver and a second implementation enumerating
activated sources agreed on all 256 seeded eight-vertex graphs.
The additive lower and witness upper bounds agreed on 253.
Counting edges gave an optimistic answer about source robustness
on three graphs. The separate seven-edge chain gave source radius
one and edge radius seven.

To test harder instances, we used 256 ten-vertex graphs with
five layers and seven possible sources. Every candidate path
requires exactly four optional steps. Two solvers again agreed
on all instances.
\begin{center}
\begin{tabular}{rrr}
\toprule
Exact source radius & Instances & Lower and upper bounds agree \\
\midrule
1 & 13 & 13 \\
2 & 157 & 117 \\
3 & 81 & 65 \\
4 & 5 & 5 \\
\bottomrule
\end{tabular}
\end{center}
The inexpensive certificate settled the exact radius in 200 of
256 cases. The two bounds differed by one on 52 cases and by
two on four. Using the edge radius four as the source radius
would have overstated robustness on 251 cases. These counts
describe the frozen graph generators and do not estimate real
log failure frequencies. No timing speedup is claimed.

The live-label solver matched direct enumeration of activated
sources on 512 additional acyclic graphs, including 76 with
no possible path. It matched enumeration on diamond chains
of one through five stages. On the 64-stage chain it returned
radius 64 using 193 states, against the analytic count of
$2^{64}$ terminal antichain entries. The latter antichain was
counted from the construction without materializing it.

\section{Choosing source audits and artifact rebuilds}
Some source groups can be audited independently. A successful audit
must establish that the group conceals no optional dependency in
the current graph. An audit that instead discovers an edge requires
updating the graph before the certificate can be used. Likewise, a
rebuild removes the artifact from suspect paths only if it excludes
all suspect inputs. Give each available audit and rebuild a cost.

Let $A$ be the selected audits and $R$ the selected rebuilds. A
candidate path with at most $k$ distinct optional labels is an
obligation. It is satisfied if its labels intersect $A$ or its
internal artifacts intersect $R$. A path with neither available
action cannot be certified away by these interventions.

\begin{theorem}
The selected audits and rebuilds rule out all hidden paths under
omission budget $k$ if and only if they satisfy every obligation.
For audit-only plans with all labels auditable, budget two reduces
to weighted vertex cover after selecting the forced singleton
labels. Even at live width zero, deciding whether the cost of a
sufficient audit plan fits a budget is NP-hard.
\end{theorem}
\begin{proof}
Remove the optional edges ruled out by the audits and the vertices
blocked by the rebuilds. If an original path with at most $k$ labels
survives, activate exactly its optional labels. This is a feasible
omission and gives a witness against restoration. Conversely, every
hidden path in the residual graph has at most $k$ labels and
survives precisely when neither intervention hit it.

With audits alone, a path with no optional labels makes the
problem infeasible. At budget two, each singleton label is forced.
Discard its satisfied obligations. Every remaining obligation has
a pair of labels, so the remaining choice is a weighted vertex
cover on the graph of these pairs. When this graph is bipartite,
a minimum cut gives the optimum. Connect the source to left
vertices and right vertices to the sink with their respective
audit costs. Connect each adjacent left and right pair with capacity
greater than the total available cost. An affordable cut chooses a
left endpoint on the sink side or a right endpoint on the source
side of each pair.

For the hardness claim, reduce ordinary vertex cover \cite{karp}.
Make a label for each vertex. For each graph edge use a separate
two-arc route from $s$ to $t$ whose two labels are its endpoints.
All such paths have two labels. No branch repeats its first label,
so every live set is empty. At budget two the audit plan must
cover exactly the original edges.
\end{proof}

An exact optimizer can find one residual witness, branch on every
eligible audit or rebuild that hits it, and retain the cheapest
valid completion. The live-label solver with predecessor pointers
can supply the witness. A witness with no eligible actions closes
that branch as infeasible. This separates the cost of finding one
counterexample from the cost of choosing the intervention set.
For an audit-only plan with at most $b$ interventions, each branch
has at most $k$ children and depth at most $b$.

As a small example, two suspect routes depend on source one and
a third depends on source two. Their distinct intermediate
artifacts cost $2,2,1$ to rebuild. The audits cost $1$ and $4$.
Auditing source one and rebuilding the third artifact costs $2$.
Using either intervention type alone costs $5$. The witness
solver matched full intervention and omission enumeration on 192
seeded graphs, and matched weighted vertex cover on 48 additional
zero-width fixtures. These are checks of the declared graph model.
A recovery team can use the mixed plan when an independent audit
is cheaper than rebuilding several artifacts. The certificate must
record what the audit established about coverage.

\newpage
\section{What a recovery drill should measure}
Record source-level outage groups and keep the label meaning
stable across tools. Inject a known complete log into a recovery
drill, withhold selected records or whole sources, and require
the team to propose a restoration before seeing the hidden
graph. Then test whether the actual hidden graph fits the
candidate family, whether the claimed budget held, and whether
the selected artifacts still receive taint. Measure unnecessary
rebuilds, time to safe service, and residual taint separately.
If common infrastructure can hide several sources together,
represent that shared failure in the budget model before using
the certificate.

\paragraph{Review status.} Independent manuscript review is pending.

\begin{thebibliography}{9}
\bibitem{broersma} Hajo Broersma, Xueliang Li, Gerhard J.
Woeginger, and Shenggui Zhang.
\emph{Paths and Cycles in Colored Graphs}.
Australasian Journal of Combinatorics 31, 299--311, 2005.
\url{https://ajc.maths.uq.edu.au/pdf/31/ajc_v31_p299.pdf}.
\bibitem{karp} Richard M. Karp.
\emph{Reducibility Among Combinatorial Problems}.
Complexity of Computer Computations, 1972.
\url{https://www.cs.umd.edu/~gasarch/BLOGPAPERS/Karp.pdf}.
\end{thebibliography}
\end{document}
